\section{Construct bounded canonical tables}
A Huffman code assigns a variable-length bit sequence to each symbol. A canonical
assignment can be reconstructed from the code lengths and symbol order; the file
does not need to store a pointer tree. Count symbols at each length, compute the
first code for each length, and place symbols in increasing length and symbol order.
Reject an assignment that asks for more codes than its prefix space can hold.

Dynamic blocks first describe a smaller alphabet used to encode the literal/length
and distance code lengths. Its repeat operators expand into the declared combined
length vector. Operator 16 needs an earlier length to repeat; operators 17 and 18
repeat zero. Check the expansion count before writing the vector, including repeats
that cross from literal/length entries into distance entries. Require the end-of-block
symbol to have a code. Bound code lengths and symbol counts before indexing tables.

Not every incomplete tree is automatically invalid: DEFLATE permits special sparse
cases, including a single distance code, and a literal-only block can have no used
distance codes. Define and test these cases rather than rejecting every positive
unused-prefix count or accepting every incomplete alphabet. The scalar smoke decoder
is a comparison implementation; its table builder is not a substitute for auditing
these format rules. Test table construction independently before combining it with
row processing.
